\documentclass{jarticle}

\usepackage{tam}
\usepackage{fleqn}
\usepackage{url}
\usepackage{bm}
\usepackage[dvipdfmx]{graphicx}
\usepackage[dvipdfmx]{color}
% \usepackage{color}
\usepackage{colordef}
\usepackage{amssymb}
\usepackage{pdfpages}
\usepackage{amsmath}

%%%%%%%%%% Command alias %%%%%%%%%%
\newcommand{\IB}{$\bullet$}
\newcommand{\IC}{$\circ$}
\newcommand{\IS}{$*$}
%%%%%%%%%%%%%%%%%%%%%%%%%%%%%%%%%%%

\title{動的計画法による電力運用最適化}
\author{神戸大学大学院 奥村侑真}
\date{}

\def\baselinestretch{1.02}
\renewcommand{\thefootnote}{\roman{footnote}}

\begin{document}
\maketitle

\section{はじめに}
図\ref{fig:tecseed}に示す電力ネットワークに対する，電力クラスタの電力運用最適化モデルを示す．

\begin{figure}[h]
 \centering
 \includegraphics[width=155mm]{../tecseed.pdf}
 \caption{テクシード構成図}
 \label{fig:tecseed}
\end{figure}

\section{準備}

基本要素およびそのスペックを表す定数(\IB を付して記す)と状態変数(\IC を付して表す)として以下のものを用意する．

\begin{itemize}
 \item[\IS] 期T$_k$ ($k=1,\ldots,K$)．本検証では各期の幅は$60$秒としている．
 \item[\IS] 太陽光発電機G
 \item[\IS] 消費機器D$^{\rm A}$
 \item[\IS] 消費機器D$^{\rm B}$
 \item[\IS] 消費機器D$^{\rm C}$
 \item[\IS] 固定バッテリーB$^{\rm F}$
  \begin{itemize}
   \item[\IB] 固定バッテリーの容量 $\overline{b}$
  \end{itemize}
 \item[\IS] 電力変換器(太陽光) E$^{\rm P}$
  \begin{itemize}
   \item[\IB] 変換効率 $\alpha^{\rm P}$
  \end{itemize}
 \item[\IS] 電力変換器(系統電力) E$^{\rm S}$
 \item[\IS] 電力変換器(消費機器A) E$^{\rm A}$
  \begin{itemize}
   \item[\IB] 変換効率 $\alpha^{\rm DA}$
  \end{itemize}
 \item[\IS] 電力変換器(消費機器B) E$^{\rm B}$
  \begin{itemize}
   \item[\IB] 変換効率 $\alpha^{\rm DB}$
  \end{itemize}
 \item[\IS] 電力変換器(消費機器C) E$^{\rm C}$
  \begin{itemize}
   \item[\IB] 変換効率 $\alpha^{\rm DC}$
  \end{itemize}
 \item[\IS] 電力変換器(固定バッテリー)E$^{\rm BF}$
  \begin{itemize}
   \item[\IB] 充電効率 $\alpha^{\rm FC}$
   \item[\IB] 放電効率 $\alpha^{\rm FD}$
   \item[\IB] 自己放電効率 $\alpha^{\rm FT}$
  \end{itemize}
 \item[\IS] 系統電力S
  \begin{itemize}
   \item[\IB] 系統買電単価 $p^{\rm BY}_{k}$
   \item[\IB] 単位時間あたりの最大系統買電量 $a^{\rm BY}$
   \item[\IB] 系統売電単価 $p^{\rm SL}_{k}$
 \end{itemize}
\end{itemize}

\section{定数・変数 \mbox{\normalsize\rm
~ (定数を「\IB」，変数を「\IC」を付して記すものとする)}}

 \begin{itemize}
  \item[\IB] 太陽光発電電力量(変換前) $g^{\rm P1}_{k}$
  \item[\IC] 太陽光発電電力量(変換後) $g^{\rm P2}_{k}$
  \item[\IC] 消費電力A (変換前) $d^{\rm A1}_{k}$
  \item[\IB] 消費電力A (変換後) $d^{\rm A2}_{k}$
  \item[\IC] 消費電力B (変換前) $d^{\rm B1}_{k}$
  \item[\IB] 消費電力B (変換後) $d^{\rm B2}_{k}$
  \item[\IC] 消費電力C (変換前) $d^{\rm C1}_{k}$
  \item[\IB] 消費電力C (変換後) $d^{\rm C2}_{k}$
  \item[\IB] 固定バッテリーのSOCの初期段階 $i_{0}$
  \item[\IB] 固定バッテリーのSOCを離散化したときの分割数$B$
  \item[\IC] 固定バッテリーの(期末の)SOC $b^{\rm F}_{k}$
  \item[\IC] 固定バッテリーの充電量(変換前) $x^{\rm FC1}_{k}$
  \item[\IC] 固定バッテリーの充電量(変換後) $x^{\rm FC2}_{k}$
  \item[\IC] 固定バッテリーの放電量(変換前) $x^{\rm FD1}_{k}$
  \item[\IC] 固定バッテリーの放電量(変換後) $x^{\rm FD2}_{k}$ \\[-2.5ex]
  \item[\IC] 系統買電力量$s^{\rm BY}_{k}$
  \item[\IC] 系統売電力量$s^{\rm SL}_{k}$
 \end{itemize}


\section{動的計画法によるアプローチ}

前節で示した定数・変数を用いて，電力運用最適化問題に対して動的計画法を用いたアプローチを検討する．

\subsection{状態空間の離散化と計算量の検討}
本モデルでは，バッテリーの充電状態(SOC)を状態変数として扱う．SOCは連続量であるが，動的計画法を適用するために離散化する必要がある．
バッテリー容量を $\overline{b} = 2742$ kWhとする．このSOCを一定の刻み幅で分割するが，分割数 $B$ (段階数)は計算精度と計算時間のトレードオフとなる．
動的計画法の計算量は，時間ステップ数を $K$，状態の段階数を $B$ とすると，一般に $O(K \cdot B^2)$ となる．
本検討では，1期を60秒とし，1週間分の最適化を考えると $K = 60 \times 24 \times 7 = 10080$ となる．$B$が大きくなると計算量が爆発的に増加するため，適切な設定が不可欠である．

例えば，最大50kWの電力を1分間（1期）充電または放電する場合，エネルギー変化量は $50 \text{kW} \times \frac{1}{60} \text{h} \approx 0.833 \text{kWh}$ となる．
この変化を十分に捉えられる精度と，現実的な計算時間を両立させる刻み幅として，ここでは0.4kWhを採用する．
これにより，SOCの段階数は
\[
B = \frac{2742 \text{ kWh}}{0.4 \text{ kWh}} + 1 = 6855 + 1 = 6856 \text{ 段階}
\]
となる．
この設定に基づくと，総計算コストは膨大になるものの，近年の計算機性能向上により，数分から数十分程度での計算が可能であると見積もられる．

\subsection{定式化}
まず，バッテリーのSOCとしてとりうる値を，0が最小値となり $\overline{b}$が最大値となるように，等間隔に $B$段階で離散化する．これにより，期 $k$ の末尾におけるSOCは以下の $B$ 個の値のいずれかをとる：
\[
0,\ \frac{1}{B-1}\overline{b},\ \frac{2}{B-1}\overline{b},\ \dots,\ \overline{b}．
\]
便宜上，期 $k$ の末尾においてSOCが $\frac{i}{B-1}\overline{b}$である状態を，「バッテリーのSOCが $i$ 段階である」と呼ぶことにする$(i=0,1,\dots ,  B - 1)$．

次に，期ごとの状態を要素に持つ集合$S_{0},S_{1},\dots,S_{K-1}$を以下のように定義する：
\[
S_k = \begin{cases}
  \left\{  s_{i} \mid i = i_{0}  \right\} ,& \quad \text{if } k = 0 \\
  \left\{ s_{i} \mid i = 0, 1, \dots, B - 1 \right\} ,& \quad \text{if } 1 \leq k < K．
\end{cases}
\]
ここで，状態$s_{i}( \in S_{k})$ は，期 $k$におけるSOCが$i$段階である状態を表す．

各状態 $ s_{i} (\in S_{k})$に対応する評価値を $F_{k}(s_{i})$，状態 $s_{j}(\in S_{k-1})$から$ s_{i} (\in S_{k})$への遷移に伴うコスト関数を$w_{k-1}(s_{j},s_{i})$とする．
このとき，コスト関数 $w_{k-1}(s_{j},s_{i})$は，期$k$における電力需給バランスから生じる売買電のコスト（または利益）として以下で与えられる：
\[
  w_{k-1}(s_{j},s_{i}) = \begin{cases}
   p_k^{\rm BY} \left[ \frac{i - j}{B-1} \cdot \overline{b} - \{ g_k^{\rm P2} - (d_k^{\rm A2} + d_k^{\rm B2} + d_k^{\rm C2}) \} \right] ，&  \text{if }  \Delta E_k \geq 0 \text{ (買電)}\\
   p_k^{\rm SL} \left[ \frac{i - j}{B-1} \cdot \overline{b} - \{ g_k^{\rm P2} - (d_k^{\rm A2} + d_k^{\rm B2} + d_k^{\rm C2}) \} \right]，& \text{otherwise (売電)}．
   \end{cases}
 \label{eq:weight}
\]
ただし，$\Delta E_k = \frac{i - j}{B-1} \cdot \overline{b} - \{ g_k^{\rm P2} - (d_k^{\rm A2} + d_k^{\rm B2} + d_k^{\rm C2}) \}$ は期$k$における系統からの正味受電量を示す．

\subsection{最適性の原理と更新式}
本問題は，最適性の原理 (Principle of Optimality) に基づいて解くことができる．最適性の原理とは，「最適政策は，その最初の状態と最初の決定がどのようであれ，最初の決定から生じる状態に関する限り，残りの決定が最適政策を構成しなければならない」というものである．

これは背理法によって証明できる．仮に，期 $k-1$ の状態 $s_j$ を経由して期 $k$ の状態 $s_i$ に至るある最適経路 $P^*$ があったとする．このとき，その部分経路である期 $k-1$ の状態 $s_j$ に至るまでの経路 $P^*_{0 \to k-1}$ が最適でなかったと仮定する．
すなわち，状態 $s_j$ に至るまでにより優れた（コストの低い）別の経路 $P'$ が存在し，その評価値を $F'_{k-1}(s_j)$ とすると，$F'_{k-1}(s_j) < F_{k-1}(s_j)$ が成り立つ．
このとき，経路 $P'$ を通り，そこから状態 $s_i$ へ遷移する新しい経路を構成すると，その総コストは $F'_{k-1}(s_j) + w_{k-1}(s_j, s_i)$ となる．これは，仮定した最適経路 $P^*$ のコスト $F_{k-1}(s_j) + w_{k-1}(s_j, s_i)$ よりも小さくなる．
これは，経路 $P^*$ が最適であるという当初の仮定と矛盾する．
したがって，最適経路の部分経路もまた最適でなければならない．

この原理に基づき，評価関数 $F_k(s_i)$ は以下の漸化式によって更新される：
\[
F_{0}(s_{i_{0}}) = 0 , \quad \text{if } k = 0
\]
\[
F_{k}(s_{i}) = \min_{s_{j} \in S_{k-1}} ( F_{k-1}(s_{j}) + w_{k-1}(s_{j},s_{i}) )  , \quad \text{if } k > 0
 \label{eq:minimize}
\]
この計算を $k=1$ から $K$ まで順次行うことで，最終的な状態における最小コストと，そこに至る最適経路（運用スケジュール）を求めることができる．

\end{document}